\section{Experiments}
\label{exp_sec}

Our experiments were conducted on \textbf{AMD EPYC 7452 32-Core Processor} with \textbf{256} GB RAM. The code was compiled using \textbf{g++ 7.5.0} without any compiler optimizations. All the times reported are in microseconds. We only show the plots here, the tables are available in the appendix \ref{sec:tables}.
\subsection{Wait-Free BFS}

We evaluate the performance of our wait-free BFS algorithm using random graphs generated with varying sizes and densities as well as real world datasets.

We evaluate two versions of the wait-free BFS algorithm, one that relies on hardware guarantees for atomicity and another that explicitly marks shared variables as atomic.

We also compare our BFS algorithm with Bader and Madduri's \cite{bader_bfs} lock-based BFS algorithm.

\subsubsection{Experimental Results}

\paragraph{Experiment 1: Varying Number of Vertices}
\label{bfs:exp1}
We fix the number of threads at 16 and vary the number of vertices in the graph keeping number of edges proportional to the number of vertices (50 times). Figure~\ref{fig:exp1} present the results.

\pgfplotstableread[col sep=space]{
    Vertices   Serial   WaitFree   WaitFreeAtomic   LockBased
    1e5        139344.6     19755.4    22067.2          61038
    2e5        311584.8     33606.8    37080            114281.8
    3e5        488414.8     47269.8    54004            169902.2
    4e5        663065.4     60732.4    70643.4          230551
    5e5        835574.8     74156.2    86906.2          284763.8
    6e5        1015243      86748.6    103275.8         338586.6
}\bfsA

\begin{figure}
    \centering
    \begin{tikzpicture}
    \begin{axis}[
        width=0.45\textwidth,
        height=0.3\textwidth,
        xlabel={Vertices},
        ylabel={Execution Time ($\mu$s)},
        legend pos=north west,
        xtick=data,
        xticklabels={1,2,3,4,5,6},
        ymajorgrids=true,
        grid style=dashed,
        enlarge x limits=0.05,
        tick label style={font=\small},
        label style={font=\small},
        legend style={font=\small},
        legend style={
            font=\scriptsize,
            cells={anchor=west},
            inner sep=2pt,
            column sep=2pt
        }
    ]
    
    \addplot table[x=Vertices, y=Serial] {\bfsA};
    \addlegendentry{Serial}
    
    \addplot table[x=Vertices, y=WaitFree] {\bfsA};
    \addlegendentry{Wait-Free}
    
    \addplot table[x=Vertices, y=WaitFreeAtomic] {\bfsA};
    \addlegendentry{Wait-Free (Atomic)}
    
    \addplot table[x=Vertices, y=LockBased] {\bfsA};
    \addlegendentry{Lock-Based}
    
    \end{axis}
    \end{tikzpicture}
    \caption{Execution time of BFS for varying number of vertices.}
    \Description{Execution time of BFS for varying number of vertices.}
    \label{fig:exp1}
\end{figure}

From Figure \ref{fig:exp1}, we observe that as the number of vertices increases, the wait-free BFS shows a consistent and substantial speedup compared to the existing parallel BFS with the wait-free version achieving up to $11\times$ speedup over the serial BFS and $4\times$ speedup over the existing parallel BFS by using $16$ threads. The version relying on hardware atomicity performs slightly better than the explicitly atomic version, but both variants achieve a significant speedup.

\paragraph{Experiment 2: Varying Graph Density}
\label{bfs:exp2}
We fix the number of vertices at $5e4$, number of threads at 64 and vary the density of the graph. Figure~\ref{fig:exp2} present the results.

\pgfplotstableread[col sep=space]{
    Density   Serial   WaitFree   WaitFreeAtomic   LockBased
    1         247092.4     20563.6    21823            66997.6
    5         1186723      52643      57167.8          195032
    10        2339696.8    89930.4    100727.2         299954.2
    25        5801718.6    192167     213233.6         506012.4
    50        11578280.2   378123.4   431672.8         814425.8
    100       23380493     737423     818978.8         1431971
}\bfsB

\begin{figure}
    \centering
    \begin{tikzpicture}
    \begin{axis}[
        width=0.45\textwidth,
        height=0.3\textwidth,
        xlabel={Density (Percent)},
        ylabel={Execution Time ($\mu$s)},
        legend pos=north west,
        xtick=data,
        xticklabels={1,5,10,25,50,100},
        ymajorgrids=true,
        grid style=dashed,
        enlarge x limits=0.05,
        tick label style={font=\small},
        label style={font=\small},
        legend style={font=\small},
        legend style={
            font=\scriptsize,
            cells={anchor=west},
            inner sep=2pt,
            column sep=2pt
        }
    ]
    
    \addplot table[x=Density, y=Serial] {\bfsB};
    \addlegendentry{Serial}
    
    \addplot table[x=Density, y=WaitFree] {\bfsB};
    \addlegendentry{Wait-Free}
    
    \addplot table[x=Density, y=WaitFreeAtomic] {\bfsB};
    \addlegendentry{Wait-Free (Atomic)}
    
    \addplot table[x=Density, y=LockBased] {\bfsB};
    \addlegendentry{Lock-Based}
    
    \end{axis}
    \end{tikzpicture}
    \caption{Execution time of BFS algorithms for varying graph densities.}
    \Description{Execution time of BFS algorithms for varying graph densities.}
    \label{fig:exp2}
\end{figure}

From Figure \ref{fig:exp2}, we observe that  the wait-free BFS maintains its efficiency as graph density increases. The speedup remains consistent across different densities, with the wait-free version offering up to $32\times$ faster execution than the sequential BFS and $2\times$ faster than existing parallel BFS using $64$ threads, especially for denser graphs. While both versions exhibit increasing overhead with higher densities, the explicitly atomic variant incurs a modest performance penalty.

\paragraph{Experiment 3: Varying Number of Threads}
\label{bfs:exp3}
We perform the experiment on a random graph with $5e5$ vertices and $2.5e7$ edges. We also perform the experiment two real world datasets, \textit{LiveJournal} \cite{livejournal} and \textit{Sina Weibo} \cite{sinaweibo}. Figure~\ref{fig:exp3_1} present the results for random graph, Figure~\ref{fig:exp3_2} present the results for \textit{LiveJournal} \cite{livejournal} graph and, Figure~\ref{fig:exp3_3} present the results for \textit{Sina Weibo} \cite{sinaweibo} graph.

\pgfplotstableread[col sep=space]{
    Threads   WaitFree   WaitFreeAtomic   LockBased
    2         471627.4     550741      2267025.2
    4         251506.2     296599.6    1036392.2
    8         136896.6     161544.2    553816.8
    16        74470.8      86648.6     285779.8
    32        48459.8      53259.2     217453.2
    64        48378.6      53296       147583.6
}\bfsC

\begin{figure}
    \centering
    \begin{tikzpicture}
    \begin{axis}[
        width=0.45\textwidth,
        height=0.3\textwidth,
        xlabel={Threads},
        ylabel={Execution Time ($\mu$s)},
        legend pos=north east,
        xtick=data,
        xticklabels={2,4,8,16,32,64},
        ymajorgrids=true,
        grid style=dashed,
        enlarge x limits=0.05,
        tick label style={font=\small},
        label style={font=\small},
        legend style={font=\small},
        legend style={
            font=\scriptsize,
            cells={anchor=west},
            inner sep=2pt,
            column sep=2pt
        }
    ]
    
    \addplot table[x=Threads, y=WaitFree] {\bfsC};
    \addlegendentry{Wait-Free}
    
    \addplot table[x=Threads, y=WaitFreeAtomic] {\bfsC};
    \addlegendentry{Wait-Free (Atomic)}
    
    \addplot table[x=Threads, y=LockBased] {\bfsC};
    \addlegendentry{Lock-Based}
    
    \end{axis}
    \end{tikzpicture}
    \caption{Execution time of BFS algorihtms for varying number of threads on random graph.}
    \Description{Execution time of BFS algorihtms for varying number of threads on random graph.}
    \label{fig:exp3_1}
\end{figure}

\pgfplotstableread[col sep=space]{
    Threads   WaitFree   WaitFreeAtomic   LockBased
    2         2126836.8     2650480.2    4365973.8
    4         1130817.4     1420545.4    2232363
    8         601362        755491.6     1192019.2
    16        362708.6      440648.8     623655.4
    32        235820.6      269026.2     399018.6
    64        172725.2      233587.4     294441.6
}\bfsD

\begin{figure}
    \centering
    \begin{tikzpicture}
    \begin{axis}[
        width=0.45\textwidth,
        height=0.3\textwidth,
        xlabel={Threads},
        ylabel={Execution Time ($\mu$s)},
        legend pos=north east,
        xtick=data,
        xticklabels={2,4,8,16,32,64},
        ymajorgrids=true,
        grid style=dashed,
        enlarge x limits=0.05,
        tick label style={font=\small},
        label style={font=\small},
        legend style={font=\small},
        legend style={
            font=\scriptsize,
            cells={anchor=west},
            inner sep=2pt,
            column sep=2pt
        }
    ]
    
    \addplot table[x=Threads, y=WaitFree] {\bfsD};
    \addlegendentry{Wait-Free}
    
    \addplot table[x=Threads, y=WaitFreeAtomic] {\bfsD};
    \addlegendentry{Wait-Free (Atomic)}
    
    \addplot table[x=Threads, y=LockBased] {\bfsD};
    \addlegendentry{Lock-Based}
    
    \end{axis}
    \end{tikzpicture}
    \caption{Execution time of BFS algorithms for varying number of threads on \textit{LiveJournal} \cite{livejournal} graph.}
    \Description{Execution time of BFS algorithms for varying number of threads on \textit{LiveJournal} \cite{livejournal} graph.}
    \label{fig:exp3_2}
\end{figure}

\pgfplotstableread[col sep=space]{
    Threads   WaitFree   WaitFreeAtomic   LockBased
    2         25707716.4     35663801    50833290.2
    4         13594550.8     19191065.4  33334874.2
    8         7132199        10601925    16136456.4
    16        5038638        7583175     10040580.6
    32        2527066.8      4045410.6   6730842.4
    64        2008005        2976349     5643021.2
}\bfsE

\begin{figure}
    \centering
    \begin{tikzpicture}
    \begin{axis}[
        width=0.45\textwidth,
        height=0.3\textwidth,
        xlabel={Threads},
        ylabel={Execution Time ($\mu$s)},
        legend pos=north east,
        xtick=data,
        xticklabels={2,4,8,16,32,64},
        ymajorgrids=true,
        grid style=dashed,
        enlarge x limits=0.05,
        tick label style={font=\small},
        label style={font=\small},
        legend style={font=\small},
        legend style={
            font=\scriptsize,
            cells={anchor=west},
            inner sep=2pt,
            column sep=2pt
        }
    ]
    
    \addplot table[x=Threads, y=WaitFree] {\bfsE};
    \addlegendentry{Wait-Free}
    
    \addplot table[x=Threads, y=WaitFreeAtomic] {\bfsE};
    \addlegendentry{Wait-Free (Atomic)}
    
    \addplot table[x=Threads, y=LockBased] {\bfsE};
    \addlegendentry{Lock-Based}
    
    \end{axis}
    \end{tikzpicture}
    \caption{Execution time of BFS for varying number of threads on \textit{Sina Weibo} \cite{sinaweibo} graph.}
    \Description{Execution time of BFS for varying number of threads on \textit{Sina Weibo} \cite{sinaweibo} graph.}
    \label{fig:exp3_3}
\end{figure}

This experiment explores the performance impact of varying thread counts on both random and real-world graphs. On both the \textit{LiveJournal} and \textit{Sina Weibo} graphs, the wait-free BFS reduces execution time by a factor of $2\times$ to $3\times$ compared to the existing parallel BFS and more than $20\times$ compared to the serial BFS on $64$ threads. The performance improvement is more pronounced in larger real-world graphs, where the algorithm's ability to distribute work efficiently across threads becomes more critical.

\paragraph{Hardware Atomicity vs. Explicit Atomic Operations}
Across all experiments, the variant that relies on hardware atomicity consistently outperforms the explicitly atomic version, albeit by a small margin. The overhead introduced by explicit atomic operations becomes more apparent in high-density graphs and larger thread counts, though the overall performance remains substantially better than the existing BFS.

\paragraph{Summary}
Our experimental evaluation demonstrates that the proposed wait-free BFS algorithm significantly outperforms the existing BFS algorithm across various graph sizes, densities, and thread counts. The wait-free BFS algorithm demonstrates robust scalability and efficiency across diverse experimental settings. It offers substantial performance improvements over existing BFS implementations, making it well-suited for large-scale graph processing in parallel environments.

\subsection{Connected Components}

We evaluate the performance of our connected components algorithms using random graphs generated with varying sizes and densities.

We also compare our algorithms with the existing connected components algorithm using wait-free parallel union-find structure proposed by Anderson and Woll \cite{anderson_dsu}.

\subsubsection{Experimental Results}

\paragraph{Experiment 1: Varying Number of Components}
\label{cc:exp1}

We fix the number of threads at 64, number of vertices at $1e6$, number of edges at $5e7$ and vary the number of components in the graph. Figure~\ref{fig:exp4} present the results.

\pgfplotstableread[col sep=space]{
    Components   Serial   Coordinated_BFS  Reduced_DSU   Existing
    1            1763774    89283    481650    494473
    10           1561431    89139    452921    473713
    50           1514085    122606   403276    487344
    100          1444857    156597   351101    456510
    500          1204890    427338   253943    472979
    1000         1165937    700199   211992    467340
}\ccA

\begin{figure}
    \centering
    \begin{tikzpicture}
    \begin{axis}[
        width=0.45\textwidth,
        height=0.3\textwidth,
        xlabel={Components},
        ylabel={Execution Time ($\mu$s)},
        legend style={at={(0.05,0.55)},anchor=west},
        xtick=data,
        xticklabels={1,10,50,100,500,1000},
        ymajorgrids=true,
        grid style=dashed,
        enlarge x limits=0.05,
        tick label style={font=\small},
        label style={font=\small},
        legend style={font=\small},
        legend style={
            font=\scriptsize,
            cells={anchor=west},
            inner sep=2pt,
            column sep=2pt
        },
        xmode=log
    ]
    
    \addplot table[x=Components, y=Serial] {\ccA};
    \addlegendentry{Serial}
    
    \addplot table[x=Components, y=Coordinated_BFS] {\ccA};
    \addlegendentry{Coordinated BFS}
    
    \addplot table[x=Components, y=Reduced_DSU] {\ccA};
    \addlegendentry{Reduced-DSU}
    
    \addplot table[x=Components, y=Existing] {\ccA};
    \addlegendentry{Existing}
    
    \end{axis}
    \end{tikzpicture}
    \caption{Execution time of connected components for varying number of components.}
    \Description{Execution time of connected components for varying number of components.}
    \label{fig:exp4}
\end{figure}

Reduced-DSU Component Merging consistently outperforms the existing algorithm across all cases, with a notable advantage when the number of components increases. Coordinated Wait-Free BFS Labeling shows excellent performance for smaller component counts but scales poorly with an increasing number of components.

\paragraph{Experiment 2: Varying Number of Vertices}
\label{cc:exp2}

We fix the number of threads at 64, number of components at $100$ and vary the number of vertices in the graph keeping number of edges proportional to the number of vertices (50 times). Figure~\ref{fig:exp5} present the results.

\pgfplotstableread[col sep=space]{
    Vertices   Serial   Coordinated_BFS   Reduced_DSU   Existing
    1e4        1453219     152577   359782     451375
    2e4        3183386     244377   712586     953522
    3e4        5019183     331578   1130626    1457245
    4e4        6930431     402903   1408706    1882721
    5e4        9495439     513397   1866233    2480675
}\ccB

\begin{figure}
    \centering
    \begin{tikzpicture}
    \begin{axis}[
        width=0.45\textwidth,
        height=0.3\textwidth,
        xlabel={Vertices (Per Component)},
        ylabel={Execution Time ($\mu$s)},
        legend pos=north west,
        xtick=data,
        xticklabels={1,2,3,4,5},
        ymajorgrids=true,
        grid style=dashed,
        enlarge x limits=0.05,
        tick label style={font=\small},
        label style={font=\small},
        legend style={font=\small},
        legend style={
            font=\scriptsize,
            cells={anchor=west},
            inner sep=2pt,
            column sep=2pt
        }
    ]
    
    \addplot table[x=Vertices, y=Serial] {\ccB};
    \addlegendentry{Serial}
    
    \addplot table[x=Vertices, y=Coordinated_BFS] {\ccB};
    \addlegendentry{Coordinated BFS}
    
    \addplot table[x=Vertices, y=Reduced_DSU] {\ccB};
    \addlegendentry{Reduced-DSU}
    
    \addplot table[x=Vertices, y=Existing] {\ccB};
    \addlegendentry{Existing}
    
    \end{axis}
    \end{tikzpicture}
    \caption{Execution time of connected components for varying number of vertices.}
    \Description{Execution time of connected components for varying number of vertices.}
    \label{fig:exp5}
\end{figure}

Since the number of components is not very high, Coordinated Wait-Free BFS Labeling performs the best. This is because the time spent on barrier synchronization is minimal, and the algorithm is as efficient as the wait-free BFS. Reduced-DSU Component Merging is the second-best, with the existing algorithm following closely.

\paragraph{Experiment 3: Varying Number of Threads}
\label{cc:exp3}

We fix the number of vertices at $5e4$ per component, number of components at $100$ and number of edges at $50$ times the number of vertices. Figure~\ref{fig:exp6} present the results.

\pgfplotstableread[col sep=space]{
    Threads   Coordinated_BFS  Reduced_DSU  Existing
    2         5771680    14365772   20918563
    4         3254026    10035562   11630302
    8         1829292    7123372    6619472
    16        991526     4880461    5024664
    32        545344     2767063    3205408
    64        490756     1814275    2728337
}\ccC

\begin{figure}
    \centering
    \begin{tikzpicture}
    \begin{axis}[
        width=0.45\textwidth,
        height=0.3\textwidth,
        xlabel={Threads},
        ylabel={Execution Time ($\mu$s)},
        legend pos=north east,
        xtick=data,
        xticklabels={2,4,8,16,32,64},
        ymajorgrids=true,
        grid style=dashed,
        enlarge x limits=0.05,
        tick label style={font=\small},
        label style={font=\small},
        legend style={font=\small},
        legend style={
            font=\scriptsize,
            cells={anchor=west},
            inner sep=2pt,
            column sep=2pt
        }
    ]
    
    \addplot table[x=Threads, y=Coordinated_BFS] {\ccC};
    \addlegendentry{Coordinated BFS}
    
    \addplot table[x=Threads, y=Reduced_DSU] {\ccC};
    \addlegendentry{Reduced-DSU}
    
    \addplot table[x=Threads, y=Existing] {\ccC};
    \addlegendentry{Existing}
    
    \end{axis}
    \end{tikzpicture}
    \vspace{-1mm}
    \caption{Execution time of connected components for varying number of threads.}
    \Description{Execution time of connected components for varying number of threads.}
    \label{fig:exp6}
\end{figure}

Similar to experiment 2, Coordinated Wait-Free BFS Labeling consistently outperforms the other algorithms across different thread counts, followed by Reduced-DSU Component Merging. The existing algorithm performs the worst.

\paragraph{Experiment 4: Varying Graph Density}
\label{cc:exp4}

We fix the number of threads at 64 and number of vertices at $1e6$. We performed the experiments on low density graphs, since large densities would result in a single connected component. Figure~\ref{fig:exp7} present the results.

\pgfplotstableread[col sep=space]{
    Density   Serial   Coordinated_BFS  Reduced_DSU  Existing
    0.0002    270152     136988494     71176     111844
    0.0004    365726     57024239      70460     169343
    0.001     391727     6325528       75408     206833
    0.002     455994     281690        88115     253037
    0.004     567360     42533         124377    309005
    0.01      985679     60102         233327    383737
}\ccD

\begin{figure}
    \centering
    \begin{tikzpicture}
    \begin{axis}[
        width=0.45\textwidth,
        height=0.3\textwidth,
        xlabel={Density (Percent)},
        ylabel={Execution Time ($\mu$s)},
        legend pos=north east,
        xtick=data,
        xticklabels={0.0002,0.0004,0.001,0.002,0.004,0.01},
        ymajorgrids=true,
        grid style=dashed,
        enlarge x limits=0.05,
        tick label style={font=\small},
        label style={font=\small},
        legend style={font=\small},
        legend style={
            font=\scriptsize,
            cells={anchor=west},
            inner sep=2pt,
            column sep=2pt
        },
        ymax=2000000,
        xmode=log
    ]
    
    \addplot table[x=Density, y=Serial] {\ccD};
    \addlegendentry{Serial}
    
    \addplot table[x=Density, y=Coordinated_BFS] {\ccD};
    \addlegendentry{Coordinated BFS}
    
    \addplot table[x=Density, y=Reduced_DSU] {\ccD};
    \addlegendentry{Reduced-DSU}
    
    \addplot table[x=Density, y=Existing] {\ccD};
    \addlegendentry{Existing}
    
    \end{axis}
    \end{tikzpicture}
    \caption{Execution time of connected components for varying graph densities.}
    \Description{Execution time of connected components for varying graph densities.}
    \label{fig:exp7}
\end{figure}

For extremely sparse graphs, Coordinated Wait-Free BFS Labeling performs the worst since the number of components is very high and time spent on barrier synchronization is significant. Reduced-DSU Component Merging is the best, followed by the existing algorithm.

\paragraph{Summary} In conclusion, for small to medium numbers of components, Coordinated Wait-Free BFS Labeling is the most efficient. However, as the number of components increases, Reduced-DSU Component Merging becomes the best choice.
